\begin{textsample}{POS Dim 1 – 2016 – Score 23.00 – 2016/t001399.txt}  \label{ex:f1_pos_051}
I was excited to be a part of \textbf{the}"Dream"\textbf{theme}, and then I found out I ' m leading off \textbf{the}"Nightmare?"section of it.
And certainly there are things about \textbf{the} climate crisis that qualify.
And I have some bad news, but I have a lot more good news.
I ' m going to propose three questions and \textbf{the} answer to \textbf{the} first one necessarily involves a little bad news.
But — hang on, because \textbf{the} answers to \textbf{the} second and third questions really are very positive.
So \textbf{the} first question is,"Do we really have to change?"And of course, \textbf{the} Apollo Mission, among other things changed \textbf{the} environmental movement, really launched \textbf{the} modern environmental movement. 18 months after this Earthrise picture was first seen on earth, \textbf{the} first Earth Day was organized.
And we learned a lot about ourselves looking back at our \textbf{planet} from space.
And one of \textbf{the} things that we learned confirmed what \textbf{the} scientists have long told us.
One of \textbf{the} most essential facts about \textbf{the} climate crisis has to do with \textbf{the} sky.
As this picture illustrates, \textbf{the} sky is not \textbf{the} vast and limitless expanse that appears when we look up from \textbf{the} ground.
It is a very thin shell of atmosphere surrounding \textbf{the} \textbf{planet}.
That right now is \textbf{the} open sewer for our industrial civilization as it ' s currently organized.
We are spewing 110 million tons of heat-trapping global warming pollution into it every 24 hours, free of charge, go ahead.
And there are many sources of \textbf{the} \textbf{greenhouse} gases, I ' m certainly not going to go through them all.
I ' m going to focus on \textbf{the} main one, but \textbf{agriculture} is involved, diet is involved, population is involved.
Management of forests, \textbf{transportation}, \textbf{the} oceans, \textbf{the} melting of \textbf{the} permafrost.
But I ' m going to focus on \textbf{the} heart of \textbf{the} problem, which is \textbf{the} fact that we still rely on dirty, carbon-based fuels for 85 percent of all \textbf{the} energy that our world burns every year.
And you can see from this image that after World War II, \textbf{the} \textbf{emission} rates started really accelerating.
And \textbf{the} accumulated amount of man- made, global warming pollution that is up in \textbf{the} atmosphere now traps as much extra heat energy as would be released by 400, 000 Hiroshima-class atomic bombs exploding every 24 hours, 365 days a year.
Fact-checked over and over again, conservative, it ' s \textbf{the} truth.
Now it ' s a big \textbf{planet}, but — that is a lot of energy, particularly when you multiply it 400, 000 times per day.
And all that extra heat energy is heating up \textbf{the} atmosphere, \textbf{the} whole earth system.
Let ' s look at \textbf{the} atmosphere.
This is a depiction of what we used to think of as \textbf{the} normal distribution of temperatures. \textbf{The} white represents normal temperature days ; 1951-1980 are arbitrarily chosen. \textbf{The} \textbf{blue} are cooler than average days, \textbf{the} red are warmer than average days.
But \textbf{the} entire curve has moved to \textbf{the} right in \textbf{the} 1980s.
And you ' ll see in \textbf{the} lower \textbf{right-hand} corner \textbf{the} appearance of statistically significant numbers of extremely hot days.
In \textbf{the} 90s, \textbf{the} curve shifted further.
And in \textbf{the} last 10 years, you see \textbf{the} extremely hot days are now more numerous than \textbf{the} cooler than average days.
In fact, they are 150 times more common on \textbf{the} \textbf{surface} of \textbf{the} earth than they were just 30 years ago.
So we ' re having record-breaking temperatures.
Fourteen of \textbf{the} 15 of \textbf{the} hottest years ever measured with instruments have been in this young century. \textbf{The} hottest of all was last year.
Last month was \textbf{the} 371st month in a row warmer than \textbf{the} 20th-century average.
And for \textbf{the} first time, not only \textbf{the} warmest January, but for \textbf{the} first time, it was more than two degrees Fahrenheit warmer than \textbf{the} average.
These higher temperatures are having an effect on animals, plants, people, ecosystems.
But on a global basis, 93 percent of all \textbf{the} extra heat energy is trapped in \textbf{the} oceans.
And \textbf{the} scientists can measure \textbf{the} heat buildup much more precisely now at all depths : deep, mid-ocean, \textbf{the} first few hundred meters.
And this, too, is accelerating.
It goes back more than a century.
And more than half of \textbf{the} increase has been in \textbf{the} last 19 years.
This has consequences. \textbf{The} first order of consequence : \textbf{the} ocean-based storms get stronger.
Super Typhoon Haiyan went over areas of \textbf{the} Pacific five and a half degrees Fahrenheit warmer than normal before it slammed into Tacloban, as \textbf{the} most destructive storm ever to make landfall.
Pope Francis, who has made such a difference to this whole issue, visited Tacloban right after that.
Superstorm Sandy went over areas of \textbf{the} Atlantic nine degrees warmer than normal before slamming into New York and New Jersey. \textbf{The} second order of consequences are affecting all of us right now. \textbf{The} warmer oceans are evaporating much more water vapor into \textbf{the} skies.
Average humidity worldwide has gone up four percent.
And it creates these atmospheric rivers. \textbf{The} Brazilian scientists call them"flying rivers."And they funnel all of that extra water vapor over \textbf{the} land where storm conditions trigger these massive record-breaking downpours.
This is from Montana.
Take a look at this storm last August.
As it moves over Tucson, Arizona.
It literally splashes off \textbf{the} city.
These downpours are really unusual.
Last July in Houston, Texas, it rained for two days, 162 billion gallons.
That represents more than two days of \textbf{the} full flow of Niagara Falls in \textbf{the} middle of \textbf{the} city, which was, of course, paralyzed.
These record downpours are creating historic floods and mudslides.
This one is from Chile last year.
And you ' ll see that warehouse going by.
There are \textbf{oil} tankers cars going by.
This is from Spain last September, you could call this \textbf{the} running of \textbf{the} cars and trucks, I guess.
Every night on \textbf{the} TV news now is like a nature hike through \textbf{the} Book of Revelation.
I mean, really. \textbf{The} insurance industry has certainly noticed, \textbf{the} losses have been mounting up.
They ' re not under any illusions about what ' s happening.
And \textbf{the} causality requires a moment of discussion.
We ' re used to thinking of linear cause and linear effect — one cause, one effect.
This is systemic causation.
As \textbf{the} great Kevin Trenberth says,"All storms are different now.
There ' s so much extra energy in \textbf{the} atmosphere, there ' s so much extra water vapor.
Every storm is different now."So, \textbf{the} same extra heat pulls \textbf{the} soil moisture out of \textbf{the} ground and causes these deeper, longer, more pervasive \textbf{droughts} and many of them are underway right now.
It dries out \textbf{the} vegetation and causes more fires in \textbf{the} western part of North America.
There ' s certainly been evidence of that, a lot of them.
More lightning, as \textbf{the} heat energy builds up, there ' s a considerable amount of additional lightning also.
These climate-related disasters also have geopolitical consequences and create instability. \textbf{The} climate-related historic \textbf{drought} that started in Syria in 2006 destroyed 60 percent of \textbf{the} farms in Syria, killed 80 percent of \textbf{the} livestock, and drove 1. 5 million climate refugees into \textbf{the} cities of Syria, where they collided with another 1. 5 million refugees from \textbf{the} Iraq War.
And along with other factors, that opened \textbf{the} gates of Hell that people are trying to close now. \textbf{The} US Defense Department has long warned of consequences from \textbf{the} climate crisis, including refugees, food and water shortages and pandemic disease.
Right now we ' re seeing microbial diseases from \textbf{the} tropics spread to \textbf{the} higher latitudes ; \textbf{the} \textbf{transportation} revolution has had a lot to do with this.
But \textbf{the} changing conditions change \textbf{the} latitudes and \textbf{the} areas where these microbial diseases can become endemic and change \textbf{the} range of \textbf{the} vectors, like mosquitoes and ticks that carry them. \textbf{The} Zika epidemic now — we ' re better positioned in North America because it ' s still a little too cool and we have a better public health system.
But when women in some regions of South and Central America are advised not to get pregnant for two years — that ' s something new, that ought to get our attention. \textbf{The} Lancet, one of \textbf{the} two greatest medical journals in \textbf{the} world, last summer labeled this a medical emergency now.
And there are many factors because of it.
This is also connected to \textbf{the} extinction crisis.
We ' re in danger of losing 50 percent of all \textbf{the} living species on earth by \textbf{the} end of this century.
And already, land-based plants and animals are now moving towards \textbf{the} poles at an average rate of 15 feet per day.
Speaking of \textbf{the} North Pole, last December 29, \textbf{the} same storm that caused historic flooding in \textbf{the} American Midwest, raised temperatures at \textbf{the} North Pole 50 degrees Fahrenheit warmer than normal, causing \textbf{the} thawing of \textbf{the} North Pole in \textbf{the} middle of \textbf{the} long, dark, winter, polar night.
And when \textbf{the} land- based ice of \textbf{the} Arctic melts, it raises \textbf{sea} level.
Paul Nicklen ' s beautiful photograph from Svalbard illustrates this.
It ' s more dangerous coming off Greenland and particularly, Antarctica. \textbf{The} 10 largest risk cities for sea-level rise by population are mostly in South and Southeast Asia.
When you measure it by assets at risk, number one is Miami : three and a half trillion dollars at risk.
Number three : New York and Newark.
I was in Miami last fall during \textbf{the} supermoon, one of \textbf{the} highest high-tide days.
And there were \textbf{fish} from \textbf{the} ocean swimming in some of \textbf{the} streets of Miami Beach and Fort Lauderdale and Del Rey.
And this happens regularly during \textbf{the} highest- tide tides now.
Not with rain — they call it"sunny-day flooding."It comes up through \textbf{the} storm sewers.
And \textbf{the} Mayor of Miami speaks for many when he says it is long past time this can be viewed through a partisan lens.
This is a crisis that ' s getting worse day by day.
We have to move beyond partisanship.
And I want to take a moment to honor these House Republicans — who had \textbf{the} courage last fall to step out and take a political risk, by telling \textbf{the} truth about \textbf{the} climate crisis.
So \textbf{the} cost of \textbf{the} climate crisis is mounting up, there are many of these aspects I haven ' t even mentioned.
It ' s an enormous burden.
I ' ll mention just one more, because \textbf{the} World Economic Forum last month in Davos, after their annual survey of 750 economists, said \textbf{the} climate crisis is now \textbf{the} number one risk to \textbf{the} global economy.
So you get central bankers like Mark Carney, \textbf{the} head of \textbf{the} UK Central Bank, saying \textbf{the} vast majority of \textbf{the} \textbf{carbon} reserves are unburnable.
Subprime \textbf{carbon}.
I ' m not going to remind you what happened with subprime mortgages, but it ' s \textbf{the} same thing.
If you look at all of \textbf{the} \textbf{carbon} fuels that were burned since \textbf{the} beginning of \textbf{the} industrial revolution, this is \textbf{the} quantity burned in \textbf{the} last 16 years.
Here are all \textbf{the} ones that are proven and left on \textbf{the} books, 28 trillion dollars. \textbf{The} International Energy Agency says only this amount can be burned.
So \textbf{the} rest, 22 trillion dollars — unburnable.
Risk to \textbf{the} global economy.
That ' s why divestment movement makes practical sense and is not just a moral imperative.
So \textbf{the} answer to \textbf{the} first question,"Must we change?"is yes, we have to change.
Second question,"Can we change?"This is \textbf{the} exciting news! \textbf{The} best projections in \textbf{the} world 16 years ago were that by 2010, \textbf{the} world would be able to install 30 gigawatts of wind capacity.
We beat that mark by 14 and a half times over.
We see an exponential curve for wind installations now.
We see \textbf{the} cost coming down dramatically.
Some countries — take Germany, an industrial powerhouse with a climate not that different from Vancouver ' s, by \textbf{the} way — one day last December, got 81 percent of all its energy from \textbf{renewable} resources, mainly solar and wind.
A lot of countries are getting more than half on an average basis.
More good news : energy storage, from batteries particularly, is now beginning to take off because \textbf{the} cost has been coming down very dramatically to solve \textbf{the} intermittency problem.
With solar, \textbf{the} news is even more exciting! \textbf{The} best projections 14 years ago were that we would install one gigawatt per year by 2010.
When 2010 came around, we beat that mark by 17 times over.
Last year, we beat it by 58 times over.
This year, we ' re on track to beat it 68 times over.
We ' re going to win this.
We are going to prevail. \textbf{The} exponential curve on solar is even steeper and more dramatic.
When I came to this stage 10 years ago, this is where it was.
We have seen a revolutionary breakthrough in \textbf{the} emergence of these exponential curves.
And \textbf{the} cost has come down 10 percent per year for 30 years.
And it ' s continuing to come down.
Now, \textbf{the} business community has certainly noticed this, because it ' s crossing \textbf{the} grid parity point.
Cheaper solar penetration rates are beginning to rise.
Grid parity is understood as that line, that threshold, below which \textbf{renewable} electricity is cheaper than electricity from burning \textbf{fossil} fuels.
That threshold is a little bit like \textbf{the} difference between 32 degrees Fahrenheit and 33 degrees Fahrenheit, or zero and one Celsius.
It ' s a difference of more than one degree, it ' s \textbf{the} difference between ice and water.
And it ' s \textbf{the} difference between markets that are frozen up, and liquid flows of capital into new opportunities for investment.
This is \textbf{the} biggest new business opportunity in \textbf{the} history of \textbf{the} world, and \textbf{two-thirds} of it is in \textbf{the} private sector.
We are seeing an \textbf{explosion} of new investment.
Starting in 2010, investments globally in \textbf{renewable} electricity generation surpassed \textbf{fossils}. \textbf{The} gap has been growing ever since. \textbf{The} projections for \textbf{the} future are even more dramatic, even though \textbf{fossil} energy is now still subsidized at a rate 40 times larger than renewables.
And by \textbf{the} way, if you add \textbf{the} projections for nuclear on here, particularly if you assume that \textbf{the} work many are doing to try to break through to safer and more acceptable, more affordable forms of nuclear, this could change even more dramatically.
So is there any precedent for such a rapid adoption of a new technology?
Well, there are many, but let ' s look at cell phones.
In 1980, AT\&T, then Ma Bell, commissioned McKinsey to do a global market survey of those clunky new mobile phones that appeared then."How many can we sell by \textbf{the} year 2000?"they asked.
McKinsey came back and said,"900, 000."And sure enough, when \textbf{the} year 2000 arrived, they did sell 900, 000 — in \textbf{the} first three days.
And for \textbf{the} balance of \textbf{the} year, they sold 120 times more.
And now there are more cell connections than there are people in \textbf{the} world.
So, why were they not only wrong, but way wrong?
I ' ve asked that question myself,"Why?"And I think \textbf{the} answer is in three parts.
First, \textbf{the} cost came down much faster than anybody expected, even as \textbf{the} quality went up.
And low-income countries, places that did not have a landline grid — they leap-frogged to \textbf{the} new technology. \textbf{The} big expansion has been in \textbf{the} developing counties.
So what about \textbf{the} electricity grids in \textbf{the} developing world?
Well, not so hot.
And in many areas, they don ' t exist.
There are more people without any electricity at all in India than \textbf{the} entire population of \textbf{the} United States of America.
So now we ' re getting this : solar panels on grass huts and new business models that make it affordable.
Muhammad Yunus financed this one in Bangladesh with micro- credit.
This is a village market.
Bangladesh is now \textbf{the} fastest-deploying country in \textbf{the} world : two systems per minute on average, night and day.
And we have all we need : enough energy from \textbf{the} Sun comes to \textbf{the} earth every hour to supply \textbf{the} full world ' s energy needs for an entire year.
It ' s actually a little bit less than an hour.
So \textbf{the} answer to \textbf{the} second question,"Can we change?"is clearly"Yes."And it ' s an ever-firmer"yes."Last question,"Will we change?"Paris really was a breakthrough, some of \textbf{the} provisions are binding and \textbf{the} regular reviews will matter a lot.
But nations aren ' t waiting, they ' re going ahead.
China has already announced that starting next year, they ' re adopting a nationwide cap and trade system.
They will likely link up with \textbf{the} European Union. \textbf{The} United States has already been changing.
All of these \textbf{coal} plants were proposed in \textbf{the} next 10 years and canceled.
All of these existing \textbf{coal} plants were retired.
All of these \textbf{coal} plants have had their retirement announced.
All of them — canceled.
We are moving forward.
Last year — if you look at all of \textbf{the} investment in new electricity generation in \textbf{the} United States, almost three-quarters was from \textbf{renewable} energy, mostly wind and solar.
We are solving this crisis. \textbf{The} only question is : how long will it take to get there?
So, it matters that a lot of people are organizing to insist on this change.
Almost 400, 000 people marched in New York City before \textbf{the} UN special session on this.
Many thousands, tens of thousands, marched in cities around \textbf{the} world.
And so, I am extremely optimistic.
As I said before, we are going to win this.
I ' ll finish with this story.
When I was 13 years old, I heard that proposal by President Kennedy to land a person on \textbf{the} Moon and bring him back safely in 10 years.
And I heard adults of that day and time say,"That ' s reckless, expensive, may well fail."But eight years and two months later, in \textbf{the} moment that Neil Armstrong set foot on \textbf{the} Moon, there was great cheer that went up in NASA ' s mission control in Houston.
Here ' s a little-known fact about that : \textbf{the} average age of \textbf{the} systems engineers, \textbf{the} controllers in \textbf{the} room that day, was 26, which means, among other things, their age, when they heard that challenge, was 18.
We now have a moral challenge that is in \textbf{the} tradition of others that we have faced.
One of \textbf{the} greatest poets of \textbf{the} last century in \textbf{the} US, Wallace Stevens, wrote a line that has stayed with me :"After \textbf{the} final ' no, ' there comes a ' yes, ' and on that ' yes ', \textbf{the} future world depends."When \textbf{the} abolitionists started their movement, they met with no after no after no.
And then came a yes. \textbf{The} Women ' s Suffrage and Women ' s Rights Movement met endless no ' s, until finally, there was a yes. \textbf{The} Civil Rights Movement, \textbf{the} movement against apartheid, and more recently, \textbf{the} movement for gay and lesbian rights here in \textbf{the} United States and elsewhere.
After \textbf{the} final"no"comes a"yes."When any great moral challenge is ultimately resolved into a binary choice between what is right and what is wrong, \textbf{the} outcome is fore-ordained because of who we are as human beings.
Ninety-nine percent of us, that is where we are now and it is why we ' re going to win this.
We have everything we need.
Some still doubt that we have \textbf{the} will to act, but I say \textbf{the} will to act is itself a \textbf{renewable} resource.
Thank you very much.
Chris Anderson : You ' ve got this incredible combination of skills.
You ' ve got this scientist mind that can understand \textbf{the} full range of issues, and \textbf{the} ability to turn it into \textbf{the} most vivid language.
No one else can do that, that ' s why you led this thing.
It was amazing to see it 10 years ago, it was amazing to see it now.
Al Gore : Well, you ' re nice to say that, Chris.
But honestly, I have a lot of really good friends in \textbf{the} scientific community who are incredibly patient and who will sit there and explain this stuff to me over and over and over again until I can get it into simple enough language that I can understand it.
And that ' s \textbf{the} key to trying to communicate.
CA : So, your talk.
First part : \textbf{terrifying}, second part : incredibly hopeful.
How do we know that all those graphs, all that progress, is enough to solve what you showed in \textbf{the} first part?
AG : I think that \textbf{the} crossing — you know, I ' ve only been in \textbf{the} business world for 15 years.
But one of \textbf{the} things I ' ve learned is that apparently it matters if a new product or service is more expensive than \textbf{the} incumbent, or cheaper than.
Turns out, it makes a difference if it ' s cheaper than.
And when it crosses that line, then a lot of things really change.
We are regularly surprised by these developments. \textbf{The} late Rudi Dornbusch, \textbf{the} great economist said,"Things take longer to happen then you think they will, and then they happen much faster than you thought they could."I really think that ' s where we are.
Some people are using \textbf{the} phrase"\textbf{The} Solar Singularity"now, meaning when it gets below \textbf{the} grid parity, unsubsidized in most places, then it ' s \textbf{the} default choice.
Now, in one of \textbf{the} presentations yesterday, \textbf{the} jitney thing, there is an effort to use regulations to slow this down.
And I just don ' t think it ' s going to work.
There ' s a woman in Atlanta, Debbie Dooley, who ' s \textbf{the} Chairman of \textbf{the} Atlanta Tea Party.
They enlisted her in this effort to put a tax on solar panels and regulations.
And she had just put solar panels on her roof and she didn ' t understand \textbf{the} request.
And so she went and formed an alliance with \textbf{the} Sierra Club and they formed a new organization called \textbf{the} Green Tea Party.
And they defeated \textbf{the} proposal.
So, finally, \textbf{the} answer to your question is, this sounds a little corny and maybe it ' s a clich{\EmojiFont �}{\EmojiFont �}, but 10 years ago — and Christiana referred to this — there are people in this audience who played an incredibly significant role in generating those exponential curves.
And it didn ' t work out economically for some of them, but it kick-started this global revolution.
And what people in this audience do now with \textbf{the} knowledge that we are going to win this.
But it matters a lot how fast we win it.
CA : Al Gore, that was incredibly powerful.
If this turns out to be \textbf{the} year, that \textbf{the} partisan thing changes, as you said, it ' s no longer a partisan issue, but you bring along people from \textbf{the} other side together, backed by science, backed by these kinds of investment opportunities, backed by reason that you win \textbf{the} day — boy, that ' s really exciting.
Thank you so much.
AG : Thank you so much for bringing me back to TED.
Thank you!

% matched lemmas: agriculture, blue, carbon, coal, drought, emission, explosion, fish, fossil, greenhouse, oil, planet, renewable, right-hand, sea, surface, terrifying, the, theme, transportation, two-third
\end{textsample}
